\fichatitulo{47 · CORPUS DE ALUCINAÇÕES}{ACL · 2024}{RAGTruth}
{\small Niu et al.. \textit{RAGTruth}. ACL · 2024. \href{https://aclanthology.org/2024.acl-long.585/}{Fonte consultada}.\par}

\campo{Problema e objetivo}
Como medir alucinações em respostas RAG de modo suficientemente granular?

\campo{Principal contribuição · síntese interpretativa}
Cria corpus anotado em nível de caso e de palavra, incorporando intensidade de alucinação.

\campo{Método / natureza da evidência}
Quase 18 mil respostas geradas por diferentes LLMs em RAG, com anotações manuais e avaliação de detectores.

\campo{Resultado ou argumento principal}
Oferece uma taxonomia operacional e exemplos de conflito entre resposta e passagem, úteis para análise de erros e treinamento de detectores.

\campo{Excertos-chave · seleção literal}
\begin{itemize}
\excerto{1}{unsupported or contradictory claims to the retrieved contents}{Resumo.}
\end{itemize}

\campo{Limitações e análise crítica}
Os domínios e idiomas do corpus não são parlamentares. As categorias precisam ser adaptadas ao regime de autoria, temporalidade e citação do Senado.

\campo{Aproveitamento na dissertação · proposta}
Metodologia: derivar categorias de erro e amostrar afirmações para avaliação humana independente.

\registro{Fonte consultada nesta preparação: PDF publicado; consulta focal do resumo, exemplo, benchmark e conclusão. Chave: \texttt{niuEtAl2024ragtruth}.}
